\documentclass[11pt]{article}
\usepackage[a4paper,margin=25mm]{geometry}
\usepackage{amsmath,amssymb,amsthm,booktabs}
\usepackage[colorlinks=true,urlcolor=blue]{hyperref}
\newtheorem{proposition}{Proposition}
\newtheorem{lemma}[proposition]{Lemma}
\newcommand{\HH}{\mathcal H}
\newcommand{\hh}{\mathcal h}
\title{Typed-state closure of the certified finite $M=5$ table}
\author{Independent local mathematical audit}
\date{8 October 2026}
\begin{document}
\maketitle

This audit concerns exactly the two Buchstab operators, the first-weighted
bank $2.4\le v\le3$, $3.2\le s_0\le5$, and upper cutoff comparison
used by \texttt{certify\_finite\_step\_table.py}. Its 24 generations
use $K=681/100$ and the separately proved initial seed $J=2327/100000$.
No new second-weighted bank is authorized. The independent numerical
replay is a separate input; numbers alone do not prove state admission.

The primitive ordinary signed-prefix interface, the first-weighted
pointwise inequality with its positive switched-product estimate, and
the four-round seed are treated as named inputs. This note checks their
state transformations and finite reuse. It does not re-prove their
moving-residue distribution theorem. The relevant sources are
\texttt{new-finite-double-sieve-recursion.tex},
\texttt{generalized-one-step-double-sieve.tex},
\texttt{second-seed-stable-prefix-extension.tex}, and
\texttt{second-seed-certified-finite-feedback.tex}.

\section{The exact class quantified over}
A state carries its marked prefix $B$, one unchanged total level $Q$,
and
\[
 T=\log_N B,\quad Q=N^\nu,\quad R=Q/B=N^L,\quad
 L=\nu-T>0,\quad \alpha=\frac{1-T}{L}\ge2.
\]
These formulas may be read as exact logarithmic identities. On a fixed
positive box their limiting parameters differ by the usual controlled
box errors. The marked slots, signs, residual atom factorials and
positive box coefficients are retained, not replaced by an arbitrary
coefficient sequence or an estimate for one individual prime tuple.

An enhanced state additionally retains all of the following data:
\begin{itemize}
\item Ordered specified slots $t_1>\cdots>t_r>0$ and every complete
square test $\sum_{i<j}t_i+2t_j<\nu$.
\item One previously licensed factorization branch, its unchanged level,
all its strict cap and gate margins, and its originally selected slots.
For a selected subproduct of exponent $U$, every old remaining-slot
test $U+\sum_{i<j}v_i+2v_j<\nu$ is retained as well.
\item A ceiling $\chi\le t_r$ below which appended primes preserve that
branch. A $W$ branch based on a future third-slot bound $c$ also has
$\chi\le c$. Once its first three actual slots exist, later primes do
not change them.
\end{itemize}
A classical state instead uses its fixed-margin total level $\nu<1/2$.
Its complete-modulus interface does not need enhanced full-well or
selected-square tests. It retains the labelled slots and the positive
lower bound on their exponents required for bounded multiplicities.

An enhanced state is admitted at physical argument $s$ only when its
cutoff $R^{1/s}$ is at most $N^\chi$, in addition to the preceding
tests. Upper arguments are at least one, lower arguments at least two.
These are conditions on the actual state, not consequences of the value
of a numerical table entry. The statement below is uniform on fixed
compact positive admitted families. It does not assert one uniform
threshold as $L$ or the prime exponents tend to zero.

\paragraph{Structural cutoff equalities.}
The physical cutoff of a Buchstab child is its newly marked prime $q$.
Consequently its new ceiling may equal $\log_Nq$; this is a structural
equality, not an entire state family that can be discarded as a thin
exception. Sifting at $q$ introduces primes strictly below $q$, so the
order-preservation argument uses ``at most the ceiling'' for the physical
cutoff and strict order for newly added primes. Complete-square and
factorization gates remain strict. Isolated integral endpoints where a
new square test becomes equality have measure zero and are approached
from the interior. They do not relax the state gates.

If a primitive seed is stated with strict comparison cutoff $L/S<\chi$,
its closed-cutoff endpoint is interpreted through the same buffered
asymptotic limit. For example, decrease $\nu$ by a positive amount
smaller than every old square and gate margin, apply the seed with its
strictly smaller comparison cutoff, and let that amount tend to zero.
The physical target count is bounded by the compared count by positivity;
the normalizing factor changes continuously by $L/(L-\varepsilon)$.
The ratio $\alpha$ increases, and all old strict gates survive for small
$\varepsilon$. This recovers the claimed limiting bound at equality.
This endpoint device does not change the levels inside an unbuffered
operator identity or reset a child's argument.

\section{Appending a marked prime preserves all retained support tests}
\begin{lemma}\label{lem:append}
Suppose an enhanced state has an added prime of exponent
$0<x\le\chi$ satisfying $2x<L$. Set
\[
 B'=Bq,\quad T'=T+x,\quad L'=L-x,\quad
 \chi'=\min(\chi,x),\quad
 \alpha'=\frac{\alpha L-x}{L-x}
           =1+\frac{\alpha-1}{1-x/L}\ge\alpha.
\]
Then the complete and selected-remaining prefix tests persist, and any
residual full-well atom below cutoff $N^{\chi'}$ at level $N^{L'}$
merges into a valid complete atom for the same retained branch.
\end{lemma}
\begin{proof}
All old specified slots precede $q$, so their tests are unchanged.
The only new specified complete-square test is $T+2x<\nu$,
equivalent to $2x<L$. For a retained selected subproduct, the selected
slots remain fixed. Before a new unselected slot, their exponent plus
all old unselected exponents equals $T$. Hence its selected-remaining
test is also $T+2x<\nu$. For later residual exponents
$y_1\ge y_2\ge\cdots$, the corresponding full or selected test is
\[
 T+x+\sum_{i<j}y_i+2y_j<\nu,
\]
which follows from the residual well-support condition at $L'=\nu-T-x$.
This checks the complete concatenated atom, not just the bare product
$Bq$. The old selected critical interval is unchanged since its selected
product and $\nu$ are unchanged. The largest prime and the retained
first three $W$ slots persist. A future-third bound persists since
$x\le\chi\le c$ and all subsequent slots are at most $x$.
\end{proof}

For a classical state, the corresponding simpler observation is
$B' d\le B'(R/q)=Q$ for every residual modulus $d\le R/q$.
Appending finitely many labelled primes preserves the bounded number
of possible label assignments at each complete modulus when all prime
exponents have a fixed positive lower bound. No enhanced ordering or
well-support condition is silently imposed on the classical family.

\section{Both Buchstab operators at the common ordinary baseline six}
Let
\[
 A(s)=\frac{sF(s)}{2e^\gamma},\qquad
 a(s)=\frac{sf(s)}{2e^\gamma}.
\]
Use $a(s)=0$ below two and the positive-count upper convention $A(s)=1$
on $1\le s\le3$. No ordinary lower-sieve assertion below two is used.
For $2\le s\le5$, Buchstab begins at $R^{1/6}$ and subtracts primes
\[
 R^{1/6}\le q<R^{1/s}.
\]
Writing $x=L/(w+1)$ gives $s-1<w\le5$, apart from endpoint conventions,
and the child data are
\begin{equation}\label{eq:Bchild}
 L'=\frac{Lw}{w+1},\quad
 \alpha'=\alpha+\frac{\alpha-1}{w},\quad
 \log_N(\hbox{child physical cutoff})=x,\quad
 \frac{L'}x=w.
\end{equation}
The new prime lies below the parent's cutoff and ceiling. Since $w>1$
in the interior, $2x<L$. Lemma~\ref{lem:append} applies and the child
physical cutoff equals its new smallest specified slot. Its ratio is
at least two.

The existing positive prime-prefix normalization gives measure
$L\,dx/[x(L-x)]$. The exact substitution in \eqref{eq:Bchild} turns
its absolute value into $dw/w$. The ordinary delay identities give
\begin{align*}
 A(6)-\int_{s-1}^{5}a(w)\frac{dw}{w}&=A(s),\\
 a(6)-\int_{s-1}^{5}A(w)\frac{dw}{w}&=a(s).
\end{align*}
Thus the ordinary baseline is six, not five; no positive correction at
six is presumed. Its cutoff $L/6$ is below the admitted physical cutoff,
so its ordinary upper and lower atoms have the same retained branch.

An upper parent uses lower child estimates in the subtracted sum and
therefore gains
\[
 \mathcal B_H[h](s)=\int_{\max(2,s-1)}^5h(w)\frac{dw}{w}.
\]
A lower parent uses upper child estimates and gains
\[
 \mathcal B_h[H](s)=\int_{s-1}^5H(w)\frac{dw}{w}.
\]
The part $1\le w<2$ in the latter is licensed. Indeed, at an upper
child one can compare to $v\ge2>w$ at the smaller cutoff
\[
       (R/q)^{1/v}\le(R/q)^{1/w}=q.
\]
Any prime $q'$ then introduced by that comparison has
$(q')^2\le R/q$, and the complete next test is
$q(q')^2\le R$. It is not the stronger, unwarranted condition $q^3\le R$.
All actual comparisons used in the finite bank have $v\le3$, so the
ordinary shift $A(w)-A(v)$ is zero in this below-two use.

For a closed target cell $[a,b]\subset[2,5]$, replace the lower limit
by $\max(2,b-1)$ or $b-1$, respectively. For every $s\in[a,b]$ the
resulting tail is a subset of its licensed child interval, so the
nonnegative integral is a uniform lower bound. The same admission
test is imposed at the actual $s$, not merely at the cell midpoint.

\section{The actual first-weighted bank and its lower children}
Let $v\in\{2.4,2.5,\ldots,3\}$ and
$s_0\in\{3.2,3.4,\ldots,5\}$. At a state admitted at $v$, the
first-weighted inequality has same-prefix upper term at argument
$s_0$ and negative singleton terms
\[
 q=R^u,\qquad 1/s_0<u<1/v.
\]
Here $u<1/v\le5/12<1/2$, so $2\log_Nq<L$ strictly, and
$q<R^{1/v}\le N^\chi$. For the lower singleton child the exact data are
\begin{equation}\label{eq:Wchild}
 L_u=L(1-u),\quad \alpha_u=\frac{\alpha-u}{1-u}\ge\alpha,\quad
 \rho=s_0(1-u),\quad
 \log_N(\hbox{actual cutoff})=\frac L{s_0}\le Lu.
\end{equation}
It is an admitted lower state whenever $\rho\ge2$. Its residual
argument is $\rho$, not $s_0$, an old-level parameter, or two. If
$\rho<2$, the construction uses count nonnegativity and zero defect.
The unchanged-prefix $H(s_0)$ term is admitted because $L/s_0<L/v$.

In particular it is \emph{not} necessary that the lower state in
\eqref{eq:Wchild} itself admits the initial seed comparison at $S=11/5$.
For its own next Buchstab child indexed by $w\ge\rho-1$,
\[
 x'=\frac{L(1-u)}{w+1}\le\frac{L(1-u)}\rho
      =\frac L{s_0}\le Lu.
\]
This preserves the old and newly marked slots. When that new child
uses the initial seed at $w\le S$, its compared cutoff has exponent
\[
             \frac{L(1-u)w}{S(w+1)}\le x'.
\]
When it uses a later upper-bank comparison $v'\ge w$, the same
calculation with $v'$ gives a cutoff at most $x'$. These are the actual
child-state admission tests needed by the induction.

As a concrete check of this distinction, take a one-prime largest-slot
state with $T=1/5$, $\nu=11/20$, $L=7/20$, $\chi=1/5$ and
$\alpha=16/7$. It is strictly below its largest-prime cap and passes
its complete square test. The bank choice $v=12/5$, $s_0=5$ and
$u=21/100$ has
\[
 x=147/2000,\quad L_u=553/2000,\quad
 \rho=79/20,\quad L/s_0=7/100<x.
\]
The lower child is admitted, but $L_u/S=553/4400>x$.
Forcing the seed comparison at this lower parent would incorrectly
reject a legal state. The next-child calculation above gives the
correct admission instead. It does not apply the seed where its
comparison test fails.

The positive switched triple is not inserted into this non-switched
typed family. It uses the separately assumed positive finite-product
theorem. Its actual ratio remains $\alpha\ge2$, and its factors have
fixed positive cutoffs. For this bank
\[
 \frac{\alpha-u_1-u_2-u_3}{u_2}\ge2v-3\ge\frac95>\frac1C,
 \qquad C=\frac{70893}{125000}.
\]
The audited Buchstab envelope therefore supplies the constant $C$.
The unit-quotient gap is at least $\alpha-3/v\ge3/4>0$.
This yields the already specified positive triple cost and no child
feedback. The resulting upper candidate is precisely
\[
 A(v)-A(s_0)+H_n(s_0)-\frac K8\overline I(v,s_0)
 +\frac12\int_{1/s_0}^{1/v}
      \frac{a(s_0(1-u))+h_n(s_0(1-u))}{u(1-u)}\,du.
\]
Here $\rho\le s_0-1\le4$, which is the range used by the ordinary
logarithmic kernel. The feedback retains its actual $\rho$.

\section{Upper comparison, the initial seed, and finite induction}
For the same prefix and the same $Q$, a proved upper bound at $v\ge s$
gives $\Phi_s\le\Phi_v$ because $R^{1/v}\le R^{1/s}$. The resulting
defect is $A(s)-A(v)+H(v)$. In the implemented bank every comparison
has $1\le s\le v\le3$, so that shift is exactly zero. If the physical
target is admitted, the smaller comparison cutoff is automatically
admitted. This includes $s<2$; no lower node below two is introduced.
Comparisons above three would require their nonzero shift and are
outside this implementation.

For each target cell whose right endpoint is $b\le v$, the same
fixed bank candidate is valid throughout the cell, subject to actual
state admission. For an $s_0$ exactly at a grid endpoint, its left
closed cell supplies a valid uniform bound; no interpolation is used.

The initial $J=2327/100000$ is the separately audited result of four
finite second-seed rounds. That theorem is used on its stated stable
enhanced or classical families, with comparison cutoff $L/S\le\chi$
on enhanced states and all square and ratio tests retained. At an
admitted physical upper argument $1\le s\le S$, this comparison
condition follows from $L/S\le L/s\le\chi$. Hence it licenses the
initial uniform cells ending at or below $S$. All other initial
defects are zero. No direct $S$-seed is presumed for a lower parent
whose comparison cutoff is inadmissible.

\begin{proposition}\label{prop:closure}
Assume the primitive interfaces and the four-round seed on the typed
families stated above. Every finite generation of the implemented
$M=5$ table gives simultaneous uniform upper and lower defect bounds
on its admitted physical states. In particular the independently
replayed 24-generation table is licensed there. It is not licensed
on all prefix families regardless of their gates.
\end{proposition}
\begin{proof}
The seed establishes generation zero. Suppose the preceding generation
holds simultaneously on the admitted family. The two Buchstab actions
use only admitted children by \eqref{eq:Bchild} and Lemma~\ref{lem:append}.
The first-weighted actions use their old-generation unchanged-prefix
upper value and lower children \eqref{eq:Wchild}; both are admitted by
the preceding analysis. Upper comparison only lowers the cutoff and
does not change $T,L,\alpha,Q$ or the normalizing factor.

Positive integral coefficients, downward rounded numerical weights,
and maxima of competing lower bounds for the same defect preserve
the inequalities. Taking a maximum chooses among estimates; it does
not add overlapping estimates. The implementation uses $H_n,h_n$
throughout each right-hand side, never a value newly computed at
generation $n+1$. Consequently these actions prove generation $n+1$
without circularity. Closed-cell choices can be made on finite common
positive marked-box partitions, retaining all gate conditions and
applying the existing uniform boxed interfaces before their limiting
refinement.

Each new Buchstab marked prime is at least $R^{1/6}$ and leaves
$L'\ge L/2$. A first-weighted marked prime is at least $R^{1/5}$ and
leaves $L'\ge7L/12>L/2$. Thus at any fixed number of generations all
marked prime exponents and residual exponents retain a fixed positive
lower bound on compact starting boxes. The fixed finite seed proof
adds only another fixed finite depth. Its factor counts, partitions
and precision parameters can be fixed before $N\to\infty$; finitely
many errors remain $o(1)$. This proves finite closure, not convergence
or uniformity at infinite depth.
\end{proof}

\section{Required exclusions and the outcome}
A future-third ceiling cannot be forgotten. For example, the two
specified exponents $(.16,.15)$, a retained third-slot bound $c=.04$,
and level $\nu=.581$ satisfy the old complete tests and lie below
$153/256-3c/8=.58265625$. But the residual exponent is $.271$.
At target $v=5/2$ the cutoff exponent is $.1084$, which exceeds $c$.
For the valid interior first-weighted parameter $s_0=4$, $u=.39$,
an appended prime has exponent $.10569$ and gives an actual third-slot
cap $.5580225$, below $.581$. The state is correctly excluded
by $L/v>\chi$ when $\chi\le c$; a table lookup that checked only the
smallest old slot $.15$ would be invalid.

Likewise, the table does not supply an unrestricted $L_1$ defect,
does not replace failed full-square or selected-remaining gates by a
numeric argument test, and does not convert an averaged marked-prefix
statement into an estimate for each tuple. The claimed $K$ in the
seed and in the table is still $6.81$; newer constants are not imported
into their proofs. No new $W_2$ feedback beyond the initial four-round
seed is used here.

With these explicit qualifications, no closure counterexample was
found for the actions actually used. The independent numerical
certificate and this state-induction proof address different steps:
together they license the finite table on its specified family.
They do not complete the distribution foundations or certify its
integration into the entire Chen functional.
\end{document}
